\documentclass[../CLT_SYK.tex]{subfiles}
\setcounter{section}{3}

\begin{document}

\section{Proof of the main results}

\subsection{Annealed free energy and coefficient asymptotics}

To identify the deterministic coefficient $A_N(\beta)$ in
\eqref{eq:finite-N-FTW-coefficient}, we use the annealed free-energy limit
from \cite{He2026}, which holds at every fixed inverse temperature.
Convexity then gives locally uniform convergence of the first derivatives.

The following is a standard convergence theorem for convex functions; see
\cite[Theorems~10.8 and~25.7]{Rockafellar1970}.

\begin{lemma} \label{lem:convex-derivatives}
Let $I\subset\mathbb R$ be open. Suppose that $h_N:I\to\mathbb R$ are
finite, differentiable, convex functions and that $h_N\to h$ pointwise,
where $h\in C^1(I)$. Then locally uniformly on $I$, 
\[
    h_N\to h
    \qquad\text{and}\quad
    h_N'\to h'.
\]

\end{lemma}

\begin{proof}
Fix $[a,b]\subset I$. For $x\in I$ and sufficiently small $s>0$,
convexity gives
\[
 \frac{h_N(x)-h_N(x-s)}{s}
 \le h_N'(x)
 \le \frac{h_N(x+s)-h_N(x)}{s}.
\]
Taking $\liminf$ and $\limsup$ as $N\to\infty$, using pointwise
convergence, and then letting $s\downarrow0$, gives
\[
 h'(x)\le\liminf_N h_N'(x)
 \le\limsup_N h_N'(x)\le h'(x).
\]
Thus $h_N'(x)\to h'(x)$ pointwise.

Each $h_N'$ is nondecreasing. Given \(\varepsilon>0\), by uniform
continuity of \(h'\) on \([a,b]\), choose a partition
\(
a=x_0<x_1<\cdots<x_m=b
\)
such that, for every \(i=0,\ldots,m-1\),
\[
 \sup_{x,y\in[x_i,x_{i+1}]}
 \lvert h'(x)-h'(y)\rvert<\varepsilon.
\]
Pointwise convergence at the finitely many $x_i$ gives, for all
sufficiently large $N$ and every $i=0,\ldots,m$,
\[
 |h_N'(x_i)-h'(x_i)|<\varepsilon.
\]
For $x\in[x_i,x_{i+1}]$, monotonicity of $h_N'$ gives
\[
 h_N'(x_i)\le h_N'(x)\le h_N'(x_{i+1}),
\]
and hence
\[
 |h_N'(x)-h'(x)|<2\varepsilon.
\]
Therefore, $h_N'\to h'$ uniformly on $[a,b]$. 

By the mean value theorem, we also have that $h_N\to h$ uniformly
on $[a,b]$.
\end{proof}
\begin{proposition} \label{prop:annealed-free-energy}
    There exists an even convex function $\ell_q\in C^1(\mathbb R)$ such that,
    locally uniformly for $\beta\in\mathbb R$,
    \begin{align}
        \frac1N\log \mathbb E Z_N(\beta)
        &\longrightarrow \ell_q(\beta),
        \label{eq:annealed-free-energy-limit}\\
        \frac1N\partial_\beta\log \mathbb E Z_N(\beta)
        &\longrightarrow \ell_q'(\beta).
        \label{eq:annealed-derivative-limit}
    \end{align}
    Moreover, for every $\beta>0$,
    \begin{equation} \label{eq:ell-derivative-bounds}
        0<\ell_q'(\beta)\leq\frac\beta q,
        \qquad
        0\leq\frac\beta q-\ell_q'(\beta)\leq\frac{\beta^3}{3},
    \end{equation}
    and, for every $\beta\in\mathbb R$,
    \begin{equation} \label{eq:ell-quadratic-error}
        0\leq\frac{\beta^2}{2q}-\ell_q(\beta)\leq\frac{\beta^4}{12}.
    \end{equation}
    In particular, as $\beta\to0$,
    \begin{equation} \label{eq:ell-small-beta}
        \ell_q(\beta)=\frac{\beta^2}{2q}+O(\beta^4),
        \qquad
        \ell_q'(\beta)=\frac\beta q+O(|\beta|^3).
    \end{equation}
    Furthermore, $\beta\ell_q'(\beta)>0$ for every $\beta\neq0$.
\end{proposition}

\begin{proof}
    We first record the normalization used in \cite{He2026}.  Set
    \begin{equation*}
        H_N^{\rm He}:=2^{-q/2}H_N,
        \qquad
        p_{N,q}^{\rm He}(s):=\frac1N\log\E\tr e^{sH_N^{\rm He}}.
    \end{equation*}
   Note that the paper also uses the normalized trace.  Since $G$ and $-G$ have the
    same law and $H_N(-G)=-H_N(G)$, for every $\beta\in\mathbb R$,
    \begin{equation} \label{eq:He-normalization-conversion}
        \frac1N\log\E Z_N(\beta)
        =\frac1N\log\E\tr e^{|\beta|H_N}
        =p_{N,q}^{\rm He}(2^{q/2}|\beta|).
    \end{equation}
    Writing $p_q:=p_{{\rm SD},q}$ for the Schwinger--Dyson pressure in
    \cite{He2026}, \cite[Theorem~1.1 and Proposition~3.3]{He2026} give
    \begin{equation*}
        p_{N,q}^{\rm He}(s)\longrightarrow p_q(s)
        \quad (s>0),
        \qquad
        p_q\in C^1((0,\infty)),
        \qquad
        \lim_{s\downarrow0}p_q(s)=0.
    \end{equation*}
    Define
    \begin{equation*}
        \ell_q(\beta):=
        \begin{cases}
            p_q(2^{q/2}|\beta|),&\beta\neq0,\\
            0,&\beta=0.
        \end{cases}
    \end{equation*}
    Then $\ell_q$ is even and continuous on $\mathbb R$, and
    \eqref{eq:He-normalization-conversion}, together with $Z_N(0)=1$,
    proves pointwise convergence in \eqref{eq:annealed-free-energy-limit}.

    We next use the derivative formula in
    \cite[Proposition~3.3]{He2026}: for any $s>0$, 
    \begin{equation} \label{eq:He-pressure-derivative}
        p_q'(s)=\frac{s}{q}\int_0^1G_{s,q}(u)^q\,\mathrm du.
    \end{equation}
    By the positive Lehmann representation and
    \cite[Theorem~2.5]{He2026}, there is a probability measure $\nu_s$
    on $[0,\infty)$ such that
    \begin{equation*}
        G_{s,q}(u)
        =\int_0^\infty
        \frac{\cosh(x(1/2-u))}{2\cosh(x/2)}\,\nu_s(\mathrm dx),
        \qquad
        \int_0^\infty x^2\,\nu_s(\mathrm dx)=2^{2-q}s^2.
    \end{equation*}
    For $\beta>0$, put
    \begin{equation*}
        K_\beta(u):=2G_{2^{q/2}\beta,q}(u).
    \end{equation*}
    The chain rule and
    \eqref{eq:He-pressure-derivative} yield
    \begin{equation} \label{eq:ell-kernel-derivative}
        \ell_q'(\beta)
        =2^{q/2}p_q'(2^{q/2}\beta)
        =\frac{2^q\beta}{q}\int_0^1G_{2^{q/2}\beta,q}(u)^q\,\mathrm du
        =\frac\beta q\int_0^1K_\beta(u)^q\,\mathrm du.
    \end{equation}
    Moreover,
    \begin{equation} \label{eq:ell-kernel-representation}
        K_\beta(u)
        =\int_0^\infty
        \frac{\cosh(x(1/2-u))}{\cosh(x/2)}\,\nu_{2^{q/2}\beta}(\mathrm dx),
        \qquad
        \int_0^\infty x^2\,\nu_{2^{q/2}\beta}(\mathrm dx)=4\beta^2.
    \end{equation}

    For $x\geq0$ and $u\in[0,1]$, the inequality $\tanh v\leq v$
    for $v\geq0$ gives
    \begin{align*}
        0
        &\leq\log\cosh(x/2)-\log\cosh(x(1/2-u))\\
        &=\int_{x|1/2-u|}^{x/2}\tanh v\,\mathrm dv\\
        &\leq\frac12\left(\frac{x^2}{4}-x^2(1/2-u)^2\right)
        =\frac{x^2}{2}u(1-u).
    \end{align*}
   Using $1-e^{-a}\leq a$ for $a\geq0$, we obtain
    \begin{equation*}
        0\leq
        1-\frac{\cosh(x(1/2-u))}{\cosh(x/2)}
        \leq1-e^{-x^2u(1-u)/2}
        \leq\frac{x^2}{2}u(1-u).
    \end{equation*}
    Integrating with respect to $\nu_{2^{q/2}\beta}$ in
    \eqref{eq:ell-kernel-representation} therefore gives
    \begin{equation} \label{eq:ell-kernel-bounds}
        0<K_\beta(u)\leq1,
        \qquad
        0\leq1-K_\beta(u)\leq2\beta^2u(1-u).
    \end{equation}
    The first inequality and \eqref{eq:ell-kernel-derivative} imply
    $0<\ell_q'(\beta)\leq\beta/q$.  For the remaining bound, use
    \begin{equation*}
        1-a^q=(1-a)\sum_{j=0}^{q-1}a^j\leq q(1-a),
        \qquad 0\leq a\leq1,
    \end{equation*}
    to deduce that
    \begin{align*}
        0
        &\leq\frac\beta q-\ell_q'(\beta)\\
        &=\frac\beta q\int_0^1\bigl(1-K_\beta(u)^q\bigr)\,\mathrm du\\
        &\leq\beta\int_0^1\bigl(1-K_\beta(u)\bigr)\,\mathrm du\\
        &\leq2\beta^3\int_0^1u(1-u)\,\mathrm du
        =\frac{\beta^3}{3}.
    \end{align*}
    This proves \eqref{eq:ell-derivative-bounds}.  Integrating from
    $0$ to $\beta$ with  $\ell_q(0)= 0$, gives
    \begin{equation*}
        0\leq\frac{\beta^2}{2q}-\ell_q(\beta)
        =\int_0^\beta\left(\frac{s}{q}-\ell_q'(s)\right)\,\mathrm ds
        \leq\int_0^\beta\frac{s^3}{3}\,\mathrm ds
        =\frac{\beta^4}{12}.
    \end{equation*}
This proves 
    \eqref{eq:ell-quadratic-error} for $\beta>0$. Since $\ell_q$ is even as we proved,  we have the claim for all $\beta\in\mathbb R$.

    The bound \eqref{eq:ell-quadratic-error} implies
    \begin{equation*}
        \lim_{\beta\to0}
        \frac{\ell_q(\beta)-\ell_q(0)}{\beta}=0.
    \end{equation*}
    Hence $\ell_q'(0)=0$.  On $\mathbb R\setminus\{0\}$, the derivative $\ell_q'$ 
    is continuous by $p_q\in C^1((0,\infty))$ and the chain rule, and is odd
    by evenness of $\ell_q$.  By \eqref{eq:ell-derivative-bounds},
    \begin{equation*}
        |\ell_q'(\beta)|\leq\frac{|\beta|}{q},
    \end{equation*}
    so it is also continuous at zero.  Thus $\ell_q\in C^1(\mathbb R)$.
    Equations \eqref{eq:ell-derivative-bounds} and
    \eqref{eq:ell-quadratic-error} give \eqref{eq:ell-small-beta}, and
    oddness of $\ell_q'$ gives $\beta\ell_q'(\beta)>0$ for $\beta\neq0$.

    It remains to prove convexity and the two local-uniform convergences.
    Set
    \begin{equation*}
        F_N(\beta):=\E Z_N(\beta),
        \qquad
        h_N(\beta):=\frac1N\log F_N(\beta),
        \qquad
        \overline\nu_N(B):=\E\tr\mathbf1_B(H_N).
    \end{equation*}
    Then $\overline\nu_N$ is a probability measure and
    \begin{equation*}
        F_N(\beta)=\int_{\mathbb R}e^{-\beta x}\,
        \overline\nu_N(\mathrm dx).
    \end{equation*}
 In particular,
    $h_N$ is finite and twice continuously differentiable.  With the
    probability measure
    \begin{equation*}
        \overline\nu_{N,\beta}(\mathrm dx)
        :=\frac{e^{-\beta x}}{F_N(\beta)}\,
        \overline\nu_N(\mathrm dx),
    \end{equation*}
    we have
    \begin{equation*}
        h_N''(\beta)
        =\frac1N\left(
            \int_{\mathbb R}x^2\,\overline\nu_{N,\beta}(\mathrm dx)
            -\left(\int_{\mathbb R}x\,
            \overline\nu_{N,\beta}(\mathrm dx)\right)^2
        \right)\geq0.
    \end{equation*}
    Thus each $h_N$ is convex, and so is its pointwise limit $\ell_q$.
    Applying Lemma~\ref{lem:convex-derivatives} with $I=\mathbb R$ proves
    both \eqref{eq:annealed-free-energy-limit} and
    \eqref{eq:annealed-derivative-limit} locally uniformly on $\mathbb R$.
\end{proof}


By \eqref{eq:finite-N-FTW-coefficient} and
\eqref{eq:annealed-derivative-limit}, locally uniformly for
$\beta\in\mathbb R$,
\begin{equation*}
  \frac{A_N(\beta)}N\longrightarrow
  \frac\beta2\ell_q'(\beta).
\end{equation*}

\begin{proof}[Proof of Theorem~\ref{thm:partition-clt}]
Let
\[
  S_N:=\frac1{\sqrt{2M_N}}\sum_{I\in\calI_N}(g_I^2-1).
\]
The classical i.i.d.\ central limit theorem gives
$S_N\Rightarrow\xi\sim\N(0,1)$.
By Lemma~\ref{lem:exact-second-chaos},
\begin{equation}\label{eq:Y-decomposition-final}
  Y_N(\beta)-1
  =c_N(\beta)S_N+\mathcal R_N(\beta),
  \qquad
  c_N(\beta)
  :=\frac{\beta\,\partial_\beta\log \mathbb E Z_N(\beta)}{\sqrt{2M_N}},
\end{equation}
where $\mathcal R_N:=\sum_{r\geq3}P_r Y_N$ is orthogonal to $S_N$.
Since $M_N\sim N^q/q!$, \eqref{eq:annealed-derivative-limit} gives
\[
  N^{q/2-1}c_N(\beta)
  \longrightarrow
  \sqrt{\frac{q!}{2}}\,\beta\ell_q'(\beta).
\]
Proposition~\ref{prop:higher-chaos-closure} gives
$N^{q/2-1}\mathcal R_N\to0$ in $L^2$.  This proves
\eqref{eq:scaled-relative-clt}.  Orthogonality in
\eqref{eq:Y-decomposition-final} also gives
\[
  \Var(Y_N)=c_N(\beta)^2+\norm{\mathcal R_N(\beta)}_2^2.
\]
Finally, for $\beta\neq0$, the identity
$\Var(Y_N)=\Var(Z_N)/(\E Z_N)^2$ and the strict positivity
$\beta\ell_q'(\beta)>0$ give \eqref{eq:standardized-clt}.
\end{proof}

\begin{proof}[Proof of Corollary~\ref{cor:free-energy-clt}]
The case $\beta=0$ is trivial, so assume $\beta\neq0$.  Set
\[
a:=\frac q2-1,
\qquad
W_N:=\log Y_N(\beta).
\]
Then
\begin{align*}
N^a\bigl(\log Z_N-\E\log Z_N\bigr)
&=
N^a\bigl(W_N-(Y_N-1)\bigr)
+N^a(Y_N-1)
-N^a\E W_N .
\end{align*}

By Theorem~\ref{thm:partition-clt},
\[
N^a(Y_N-1)
\Longrightarrow
\sqrt{\frac{q!}{2}}\,
\beta\ell_q'(\beta)\,\xi,
\qquad
\xi\sim\mathcal N(0,1).
\]
Moreover, $Y_N-1\to0$ in probability and
$N^a(Y_N-1)$ is tight.  Since whenever $|x|\le1/2$, 
\[
|\log(1+x)-x|\le x^2,
\]
it follows that
\[
N^a\bigl(W_N-(Y_N-1)\bigr)
\xrightarrow{\mathrm p}0.
\]
Thus it remains to prove
\[
\E W_N=o(N^{-a}).
\]

Write, as in \eqref{eq:Y-decomposition-final},
\[
Y_N-1=c_NS_N+\mathcal R_N,
\]
where
\[
c_N=\frac{\beta\,\partial_\beta\log \mathbb E Z_N(\beta)}{\sqrt{2M_N}}= O(N^{-a}),\quad S_N=\frac1{\sqrt{2M_N}}\sum_{I\in\mathcal I_N}(g_I^2-1)
,\quad
\|\mathcal R_N\|_2^2=O(N^{-3a})
\]
by Proposition~\ref{prop:higher-chaos-closure}. 
We note that 
\[
\sup_N \E S_N^4<\infty.
\]  Orthogonality therefore
gives
\[
\E(Y_N-1)^2=O(N^{-2a}).
\]
On the other hand, we estimate
\begin{align*}
\mathbb P(Y_N<1/2)
&\le
\mathbb P(|c_NS_N|>1/4)
+
\mathbb P(|\mathcal R_N|>1/4)\\
&\le
C c_N^4\E S_N^4
+
C\|\mathcal R_N\|_2^2\\
&=O(N^{-3a}).
\end{align*}

Next we will prove that 
\[
Z_N(\beta;G)\ge1.
\]
Choose $x\in I$. 
The unitary $U_x$ from Lemma~\ref{lem: parity rule lemma} satisfies
\[
U_x\psi_IU_x^*=-\psi_I,
\]
and hence $\tr\psi_I=0$.  Therefore
\(
\tr H_N(G)=0.
\) 
Jensen's inequality applied to the eigenvalues of $H_N(G)$ now gives
\[
Z_N(\beta;G)
=
\tr e^{-\beta H_N(G)}
\ge
e^{-\beta\tr H_N(G)}
=1.
\]

Define
\[
h(y):=y-1-\log y,
\qquad y>0.
\]
Since $\E Y_N=1$,
\[
-\E W_N
=
\E h(Y_N).
\]
On $\{Y_N\ge1/2\}$,
\[
h(Y_N)\le (Y_N-1)^2,
\]
and therefore
\[
\E\!\left[
h(Y_N)\mathbf 1_{\{Y_N\ge1/2\}}
\right]
=O(N^{-2a}).
\]

On the other hand, if $Y_N<1/2$, then
\[
h(Y_N)\le-\log Y_N.
\]
Since $Z_N\ge1$,
\[
-\log Y_N
=
\log\E Z_N-\log Z_N
\le
\log\E Z_N.
\]
By Proposition~\ref{prop:annealed-free-energy},
\(
\log\E Z_N=O(N),
\) 
we have
\begin{align*}
\E\!\left[
h(Y_N)\mathbf 1_{\{Y_N<1/2\}}
\right]
&\le
O(N)\,\mathbb P(Y_N<1/2)=
O(N^{1-3a}).
\end{align*}
Since $q\ge4$, we have $a\ge1$, and hence
\(
1-3a\le-2a.
\) 
Consequently,
\[
0\le-\E W_N
=
O(N^{-2a})
=
O(N^{2-q})
=
o(N^{-a}).
\]

Combining the three terms above and applying Slutsky's theorem yields
\[
N^{q/2-1}
\bigl(
\log Z_N(\beta)-\E\log Z_N(\beta)
\bigr)
\Longrightarrow
\mathcal N\left(
0,\frac{q!}{2}\beta^2(\ell_q'(\beta))^2
\right).
\]
\end{proof}
\iffalse
{\color{blue} 
\begin{lemma}\footnote{SN: I think we can remove the blue.}
  Write 
  \begin{equation*}
    W_N \coloneqq \log Y_N(\beta).
  \end{equation*}
  Then, for $\beta\neq0$ and any $\lambda, t > 0$,
  \begin{align}
    \mathbb P\!\left(W_N - \E W_N < -t \right)
    &\leq \exp\left(-\frac{qt^2}{2\beta^2}\right),
    \label{eq:centered W_N tail}\\
    \E\exp\!\left(\lambda(W_N - \E W_N)\right) 
    &\leq \exp\!\left(\frac{\lambda^2 \beta^2}{2q}\right).
    \label{eq:centered W_N exp}
  \end{align}
\end{lemma}

\begin{proof}
To apply the Gaussian concentration
  inequality, we view $Z_N(\beta;G)$ as a function of
  $G := (g_I)_{I \in \calI_N}$ and establish the Lipschitz bound for $W_N$.
  Fix $I \in \calI_N$.
  By the one-variable Duhamel formula \eqref{eq:one-variable-Duhamel-formula},
  we deduce that 
  \begin{align*}
    \partial_{g_I} Z_N 
    &= 
    \tr \partial_{g_I} e^{-\beta H_N(G)}\\
    &=
    -\beta     
    \tr \!\left( \partial_{g_I} H_N(G) \cdot e^{-\beta H_N(G)} \right)\\ 
    &=
    -\beta \sigma_N \tr\bigl( e^{-\beta H_N} \psi_I \bigr).
  \end{align*}
  Write $\langle \cdot \rangle_G$ 
  for the expectation with respect to the normalized Gibbs state,
  i.e.,
  \begin{equation*}
    \langle A \rangle_G
    \coloneqq 
    \frac{\tr(A e^{-\beta H_N(G)})}{Z_N(\beta;G)}.
  \end{equation*}
  Then we obtain 
  \begin{equation} \label{eq:GC1}
    \partial_{g_I} W_N = -\beta \sigma_N \langle \psi_I \rangle_G.
  \end{equation}
  Fix an arbitrary subset $S \subseteq [N]$ with $|S| = q-1$.
  Set $A_i \coloneqq \psi_{S \cup \{i\}}$ for each $i \notin S$.
  Note that, by the relation $\{\psi_i, \psi_j\} = 2\delta_{ij} \I$,
  it holds that   
  \begin{equation} \label{eq:A_i relation}
    A_i^2 = \I,
    \qquad 
    A_i A_j + A_j A_i = 0 \quad \text{if}\quad i \neq j.
  \end{equation}
  This yields 
  \begin{equation*}
    \left(\sum_{i \notin S} \langle A_i \rangle_G A_i\right)^2
    = 
    \sum_{i \notin S}  \langle A_i \rangle_G^2 A_i^2
    +
    \sum_{\substack{i,j \notin S \\ i < j}}
    \langle A_i \rangle_G \langle A_j \rangle_G
    (A_iA_j + A_j A_i) 
    = 
    \sum_{i \notin S}  \langle A_i \rangle_G^2 \I.
  \end{equation*}
  If we write $\|\cdot\|$ for the operator norm,
  then, by \eqref{eq:A_i relation},
  we obtain 
  \begin{equation*}
    \left\|\left(\sum_{i \notin S} \langle A_i \rangle_G A_i\right)^2\right\|
    = 
    \sum_{i \notin S}  \langle A_i \rangle_G^2.
  \end{equation*}
  Since $A_i$ is self-adjoint,
  it follows that 
  \begin{equation} \label{eq:GC2}
    \left\| \sum_{i \notin S} \langle A_i \rangle_G A_i \right\| 
    = \sqrt{\sum_{i \notin S}  \langle A_i \rangle_G^2}.
  \end{equation}
  On the other hand, H\"{o}lder's inequality for Schatten norms yields 
  \begin{equation} \label{eq:GC3}
    \left\langle \sum_{i \notin S} \langle A_i \rangle_G A_i \right\rangle_G 
    \leq
    \left\| \sum_{i \notin S} \langle A_i \rangle_G A_i \right\|.
  \end{equation}
  Since the left-hand side equals $\sum_{i \notin S} \langle A_i \rangle_G^2$,
  we deduce from \eqref{eq:GC2} and \eqref{eq:GC3} that 
  \begin{equation} \label{eq:GC4}
    \sum_{i \notin S} \langle A_i \rangle^2 \leq 1.
  \end{equation}
  For each $I \in \calI_N$,
  there exist exactly $q$ subsets $S \subseteq [N]$ satisfying $I = S \cup \{i\}$.
  Hence,
  \begin{equation*}
    \sum_{\substack{S \subseteq [N]\\ |S| = q-1}} 
    \sum_{i \notin S} 
    \langle \psi_{S \cup \{i\}}\rangle_G^2 
    = 
    q \sum_{I \in \calI_N} \langle \psi_I\rangle_G^2.
  \end{equation*} 
  This, combined with \eqref{eq:GC4}, yields 
  \begin{equation*}
    \sum_{I \in \calI_N} \langle \psi_I \rangle_G^2
    \leq 
    \frac{1}{q} 
    \sum_{\substack{S \subseteq [N]\\ |S| = q-1}} 1 
    = 
    \frac{1}{q} \binom{N}{q-1}
    \leq    
    \frac{N^{q-1}}{q!}.
  \end{equation*}
  From this and \eqref{eq:GC1}, we deduce that 
  \begin{equation*}
    \left\| \nabla W_N \right\|_2^2 
    \leq
    \frac{\beta^2 \sigma_N^2 N^{q-1}}{q!}
    = 
    \frac{\beta^2}{q},
  \end{equation*}
  where the last equality follows from the definition of $\sigma_N^2$.
  This implies that $W_N(G)$ is $|\beta|/\sqrt q$-Lipschitz as a function of $G = (g_I)_{I \in \calI_N}$.
  Hence, the Gaussian concentration inequality gives the desired result.
\end{proof}

\begin{proof}[Proof of Corollary~\ref{cor:free-energy-clt}]
The case $\beta=0$ is trivial, so assume $\beta\neq0$.
Put
\[
  a:=\frac q2-1,\qquad
  W_N:=\log Y_N(\beta).
\]
It then holds that 
\begin{align*}
  \log Z_N - \E\log Z_N 
  &=
  (\log Z_N - \log \E Z_N) + (\log \E Z_N - \E \log Z_N)\\
  &=
  W_N - \E W_N\\
  &= 
  \bigl(W_N-(Y_N(\beta)-1)\bigr)
  +(Y_N(\beta)-1)-\E W_N,
\end{align*}
which yields 
\begin{equation*}
  N^a(\log Z_N - \E \log Z_N)
  =
  N^a\bigl(W_N-(Y_N(\beta)-1)\bigr)
  +N^a(Y_N(\beta)-1)-N^a \E W_N.
\end{equation*}
By Theorem~\ref{thm:partition-clt},
we have 
\begin{equation} \label{eq:cor1}
    N^a(Y_N(\beta)-1)
    \Longrightarrow
    \sqrt{\frac{q!}{2}}\,\beta\ell_q'(\beta)\,\xi.
\end{equation}
Let $C_0 \in (0,\infty)$ be such that 
$|\log (1+x) - x| \leq C_0x^2$ for all $x \in [-1/2, 1/2]$.
Since $W_N = \log Y_N(\beta)=\log(1+(Y_N(\beta)-1))$,
it follows that 
\begin{equation*}
  \E \bigl[ 1 \wedge (N^a|W_N-(Y_N(\beta)-1)|) \bigr]
  \leq
  \E \bigl[ 1 \wedge (C_0 N^a|Y_N(\beta)-1|^2) \bigr]
  + 
  \mathbb P(|Y_N(\beta)-1| > 1/2).
\end{equation*}
By \eqref{eq:cor1},
$|Y_N(\beta)-1| \xrightarrow{\mathrm{p}} 0$ and
$(N^a |Y_N(\beta)-1|)_{N \geq 1}$ is tight.
Thus, 
the right-hand side of the above inequality converges to $0$ as $N \to \infty$.
Hence, $N^a(W_N-(Y_N(\beta)-1)) \xrightarrow{\mathrm{p}} 0$.
Therefore, it suffices to show that $\E W_N = o(N^{-a})$.
% To this end,
% we first establish a Lipschitz bound for $G\mapsto W_N(G)$.  
% Thus $W_N=\log(1+(Y_N-1))$, and Theorem~\ref{thm:partition-clt} gives
% We claim that
% \begin{equation}\label{eq:log-centering-negligible}
%   \E W_N=o(N^{-(q/2-1)}).
% \end{equation}
% Differentiation gives
% \[
%   \partial_{g_I}W_N
%   =-\beta\sigma_N\langle\psi_I\rangle_{\beta,g}.
% \]
% For each $(q-1)$-set $S$, the operators
% $(\psi_{S\cup\{i\}})_{i\notin S}$ are pairwise anticommuting Hermitian
% involutions.  If $A_1,\ldots,A_m$ are any such operators and
% $a_j=\langle A_j\rangle$ in a state, set $B=\sum_j a_jA_j$ and
% $s=\sum_j a_j^2$.  Pairwise anticommutation gives
% \[
%   B^2=s\I,
%   \qquad
%   s=\langle B\rangle\leq\norm{B}=\sqrt{s},
% \]
% and hence $s\leq1$.  Applying this observation for each $S$ and noting that
% every $I\in\calI_N$ is counted exactly $q$ times yields
% \[
%   \sum_{I\in\calI_N}|\langle\psi_I\rangle_{\beta,g}|^2
%   \leq \frac1q\binom N{q-1}.
% \]
% Consequently,
% \[
%   \sum_{I\in\calI_N}|\partial_{g_I}W_N|^2
%   \leq
%   \frac{\beta^2\sigma_N^2}{q}\binom N{q-1}
%   \leq\frac{\beta^2}{q}.
% \]
% Writing $\mu_N=\E W_N$, the Gaussian concentration inequality therefore
% implies, for every $\lambda\in\mathbb R$,
% \begin{equation}\label{eq:WN-subgaussian}
%   \E e^{\lambda(W_N-\mu_N)}
%   \leq \exp\left(\frac{\lambda^2\beta^2}{2q}\right).
% \end{equation}

We first claim that 
\begin{equation} \label{eq:EWN range}
  -\frac{\beta^2}{2q}\leq \E W_N\leq0.
\end{equation}
Indeed, the upper bound follows from Jensen's inequality,
while the lower bound follows from \eqref{eq:centered W_N exp} with $\lambda=1$ and the identity $\E e^{W_N}=1$.
 By Theorem~\ref{thm:partition-clt},
\begin{equation} \label{eq:deltaN asymp}
  \delta_N^2 := \E (Y_N(\beta)-1)^2= O(N^{-(q-2)}).
\end{equation}
Henceforth, we consider sufficiently large $N$ so that $\delta_N \leq e^{-1}$.
Since $e^{W_N}-1=Y_N(\beta)-1$,
Markov's inequality implies that for any $t\geq 1$, 
\begin{equation} \label{eq:W_N tail1}
  \mathbb P(W_N<-t)
  \leq \frac{\delta_N^2}{(1-e^{-t})^2}
  \leq C \delta_N^2.
\end{equation}

On the other hand, 
by $-\E W_N \leq K_\beta:= \beta^2/(2q)$ due to \eqref{eq:EWN range} and \eqref{eq:centered W_N tail}, for any $t>K_\beta$,
\begin{equation} \label{eq:W_N tail2}
  \mathbb P(W_N<-t)
  \leq
  \exp\left(-\frac{q(t-K_\beta)^2}{2\beta^2}\right).
\end{equation}
Set
\[
  T_N:=\max\left\{
    1,\,
    K_\beta+\frac{2|\beta|}{\sqrt q}
    \sqrt{\log(1/\delta_N)}
  \right\}.
\]
We represent
\[
  \E\bigl[|W_N|\,\mathbf1_{\{W_N<-1\}}\bigr]
  =
  \mathbb P(W_N<-1)+
  \int_1^\infty\mathbb P(W_N<-t)\,\mathrm dt.
\]
Applying \eqref{eq:W_N tail1} to the first term and the integral over $[1,T_N]$
and applying \eqref{eq:W_N tail2} to the integral over $(T_N, \infty)$,
we obtain
\begin{equation*}
  \E\bigl[|W_N|\,\mathbf1_{\{W_N<-1\}}\bigr]
  \leq     
  C T_N \delta_N^2
  +
  \int_{T_N}^\infty \exp\left(-\frac{q(t-K_\beta)^2}{2\beta^2}\right) \mathrm dt.
\end{equation*}
By the Gaussian tail bound $\int_a^\infty e^{-t^2} \mathrm dt \leq C e^{-a^2}$ for any $a \geq 0$,
the second term is bounded by $C \delta_N^2$. 
Consequently,
\begin{equation}\label{eq:WN-negative-tail}
  \E\bigl[|W_N|\,\mathbf1_{\{W_N<-1\}}\bigr]
  \leq C\delta_N^2\sqrt{\log(1/\delta_N)}.
\end{equation}
Now set $h(w):=e^w-1-w$.  
Note that
\[
  h(w)\leq C(e^w-1)^2\quad (w\geq-1),
  \qquad
  h(w)\leq |w|\quad (w<-1).
\]
This, combined with $e^{W_N} - 1 = Y_N(\beta)-1$ and
\eqref{eq:WN-negative-tail}, yields
\begin{equation*}
  \E h(W_N) 
  \leq    
  C\E (Y_N(\beta)-1)^2
  +\E\bigl[|W_N|\, \mathbf1_{\{W_N<-1\}}\bigr]
  \leq
  C\delta_N^2\sqrt{\log(1/\delta_N)}.
\end{equation*}
Moreover,
since $h(W_N) = Y_N(\beta)-1-W_N$ and $\E(Y_N(\beta)-1) = 0$,
we have $-\E W_N = \E h(W_N)$.
Therefore, by \eqref{eq:EWN range},
we deduce that 
\begin{equation*}
  0 \leq -\E W_N \leq C\delta_N^2\sqrt{\log(1/\delta_N)}.
\end{equation*}
From this and \eqref{eq:deltaN asymp},
it follows that 
\[\E W_N
  =O\!\left(N^{2-q}\sqrt{\log N}\right)
  =o(N^{-(q/2-1)}),
\]
which completes the proof.
\end{proof}
}
\fi
\end{document}
